\section{选题的背景与意义}
\subsection{选题的背景}
随着数字化办公的普及，幻灯片已成为团队协作、学术汇报、项目展示等场景的核心载体，但传统PPT制作工具存在操作复杂度高、格式兼容差、不支持轻量化标记语言等问题。基于Markdown语法的Slidev框架凭借对LaTeX公式、代码高亮、Mermaid流程图的原生支持，大幅简化了幻灯片制作流程，成为轻量化幻灯片制作的主流选择。

然而当前Slidev生态缺乏面向团队的协作能力，且市面上主流AI幻灯片制作工具普遍存在多人协作困难、文本资源版本控制缺失、自定义模板支持不足等问题，无法满足组会汇报、毕设答辩等场景下的团队化制作需求，因此亟需一款适配Slidev特性的多人协作智能幻灯片平台。

\subsection{选题的意义}
本选题的现实意义主要体现在以下三个方面：
\begin{enumerate}[label=(\arabic*)]
    \item 解决Slidev框架下多人协作的空白，通过角色权限分配\cite{zhang2013rbac,wei2025rbaccloud}、版本控制与回溯功能\cite{SXDS202505026,XMGJ202006028}，满足团队场景下幻灯片协同制作的核心需求，避免传统工具中文件传输混乱、版本不一致的问题。
    \item 优化幻灯片制作效率，整合AI内容生成与格式优化能力，结合云端资源管理\cite{li2015filenet,zhang2024javaweb}实现图片、模板的统一存储与调用，降低非专业人员的制作门槛。
    \item 弥补现有AI幻灯片工具的短板，支持自定义模板适配（如组会、毕设答辩模板），兼顾易用性与个性化需求。
\end{enumerate}


\section{研究的基本内容与拟解决的主要问题}

\subsection{研究的基本内容}
本研究旨在基于Slidev框架设计并开发一款支持多人协作的智能幻灯片平台，整体目标是实现"协作制作+云端管理+AI辅助+格式兼容"的核心能力，具体拆解为以下子目标，各子目标环环相扣形成完整的平台体系：
\begin{enumerate}[label=(\arabic*)]
    \item 多人协作功能开发：实现多人实时协同编辑\cite{DNZS201028018}、角色权限分配\cite{wu2002rbacweb}、幻灯片版本控制与回溯功能。其中协同编辑需保证多端操作的一致性，权限分配需适配团队内不同角色（如编辑、查看、管理员）的操作范围，版本控制需兼顾简易性与实用性，避免Git等复杂工具的使用门槛。
    \item 云端资源管理模块开发：搭建云端存储体系，实现图片、Slidev模板等资源的统一存储、共享与快速调用，支持关键词检索功能，解决团队内资源分散、复用困难的问题。
    \item AI辅助功能集成：对接通用大模型API实现幻灯片内容生成、格式优化能力，适配Slidev的Markdown语法规范，确保AI生成内容可直接导入平台使用。
    \item Slidev格式兼容适配：深度兼容Slidev原生功能，确保平台内制作的幻灯片支持LaTeX公式渲染、代码高亮、Mermaid流程图展示，与原生Slidev项目无缝衔接\cite{GYKJ201303028}。
\end{enumerate}

\subsection{拟解决的主要问题}
针对上述研究内容，本研究需解决的核心问题与难点如下：
\begin{enumerate}[label=(\arabic*)]
    \item 多人实时协同的一致性问题：Slidev基于Markdown文件存储，多端同时编辑易出现内容冲突，需设计轻量级的冲突检测与合并机制，兼顾实时性与数据一致性，这是协作功能的核心难点。
    \item 云端资源与Slidev的适配问题：需设计符合Slidev资源引用规范的云端存储路径，确保云端图片、模板可直接被Slidev项目调用，避免格式转换或路径修改的额外操作。
    \item AI生成内容的语法适配问题：现有AI生成的文本内容难以直接适配Slidev的Markdown语法，需开发语法转换模块，将AI生成的通用文本转换为符合Slidev规范的代码块、公式块等格式。
    \item 版本控制的简易性平衡问题：需在版本记录的完整性与操作的简易性之间找到平衡，设计非技术人员可快速上手的版本回溯、对比功能，避免专业版本工具的使用复杂度。
\end{enumerate}

本研究的创新点在于：将轻量级版本控制与Slidev生态结合，提出适配非技术团队的协作方案；针对Slidev语法定制AI内容生成规则，提升AI与轻量化幻灯片工具的适配性。

\section{研究的方法与措施}

针对研究内容与核心问题，本研究采用的方法与措施如下：
\begin{enumerate}[label=(\arabic*)]
    \item 多人协作功能：采用"前端WebSocket+后端增量同步"的方案实现实时协同，前端监听编辑操作并发送增量数据，后端接收后通过文本块对比算法进行冲突检测，冲突时触发手动合并提示；角色权限分配参考成熟的权限管理设计规范\cite{kuhlmann2013uml,li2011reengineering}，设计基于角色的权限矩阵，限定不同角色的操作范围；版本控制采用基于文件快照的轻量级方案，定期保存版本快照并记录操作日志，支持快照对比与回溯。
    \item 云端资源管理：采用对象存储服务搭建云端资源库，设计统一的资源访问API，前端通过API调用资源并自动转换为Slidev可识别的相对路径；资源检索采用关键词索引机制，优化资源分类逻辑，确保资源调用效率。
    \item AI辅助功能：构建针对Slidev语法的提示词模板，引导AI生成符合Markdown规范的内容；开发语法校验与转换模块，对AI生成内容进行格式检查与修正，参考插件化系统的设计思路\cite{cai2022microfrontend}实现AI模块的可插拔扩展，便于后续功能迭代。
    \item 格式兼容适配：梳理Slidev官方语法规则并编写适配层代码，通过单元测试覆盖LaTeX公式、Mermaid流程图等核心场景，参考WEB文档协作系统的兼容性设计思路\cite{tan2012web,li2020cscw}优化渲染效果，确保与原生Slidev功能无差异。
\end{enumerate}

\section{研究的总体安排与进度}

本研究的总体安排围绕"需求分析-设计开发-测试优化-论文撰写"四个阶段展开，具体进度节点如下：
\begin{enumerate}[label=\arabic*.]
    \item 2025.12.30-2026.01.15：完成用户需求调研，输出需求规格说明书；完成平台架构、数据库及接口设计，确定核心问题解决方案并输出详细设计文档；
    \item 2026.01.16-2026.02.10：完成多人实时协同编辑、角色权限分配模块开发与单元测试；
    \item 2026.02.11-2026.03.05：完成版本控制与回溯、云端资源存储及共享模块开发与单元测试；
    \item 2026.03.06-2026.03.30：完成AI内容生成/格式优化模块开发及Slidev格式兼容适配，开展全模块联调测试；
    \item 2026.04.01-2026.04.15：开展功能、性能及用户体验测试，收集反馈并完成优化迭代；
    \item 2026.04.16-2026.04.20：输出测试报告，解决遗留问题，通过中期检查；
    \item 2026.04.21-2026.05.10：整理开发文档与测试数据，完成毕业论文初稿撰写及修改完善；
    \item 2026.05.11-2026.05.17：完成毕业论文终稿，准备答辩材料并参加毕业论文答辩；
    \item 2026.05.18-2026.05.20：完善毕设全套材料，按学校要求完成归档与最终提交。
\end{enumerate}

\section{主要参考文献}
